\subsubsection{中间构型怎样分布}

若填写$M$个中间构型，连同两端共有$M+2$个结构。以$j=0,\ldots,M+1$编号。对于非周期模型，或已选定连续周期映像的两端，线性插值的位置关系为
\begin{equation}
  \alpha_j=\frac{j}{M+1},\qquad
  \mathbf r_i^{(j)}=(1-\alpha_j)\mathbf r_i^{(0)}
     +\alpha_j\mathbf r_i^{(M+1)}.
  \label{eq:neb-image-interpolation}
\end{equation}
$i$表示同一个原子在两端对应的序号，$\alpha_j$是无量纲路径参数，位置以\AA 表示。周期结构还要按所用周期边界处理端点位移，不能将同一原子在两个不同周期映像中的坐标直接当作很长的迁移距离。

IDPP进一步使每个构型中的原子对距离接近插值目标：
\begin{equation}
\begin{aligned}
  d_{ik}^{*,(j)}&=(1-\alpha_j)d_{ik}^{(0)}
                      +\alpha_j d_{ik}^{(M+1)},\\
  E_{\mathrm{IDPP}}^{(j)}&=
    \sum_{i<k}\frac{\bigl(d_{ik}^{(j)}-d_{ik}^{*,(j)}\bigr)^2}
                         {\bigl(d_{ik}^{(j)}\bigr)^4}.
\end{aligned}
\label{eq:idpp-pair-objective}
\end{equation}
其中$i,k$为原子序号，$j$为构型序号；$d_{ik}^{(j)}$是当前原子对距离，$d_{ik}^{*,(j)}$是对应的目标距离。分母使过短的原子间距受到较强惩罚，因此IDPP往往能改善线性插值中的碰撞。上式给出单个构型的原子对目标；路径优化还会考虑相邻构型之间的联系。

$E_{\mathrm{IDPP}}$是几何插值的辅助目标，其值用于评价路径几何；反应能垒需由后续能量计算确定。窗口中的\ui{IDPP力收敛阈值}用于这一步辅助优化，与电子结构计算中的原子力阈值分别设置。$M=5$时，路径参数依次为$0,1/6,2/6,\ldots,1$；本例的00至06正好对应七个结构。
